\documentclass[10pt,a4paper,withhyper,normalphoto]{altacv}
\geometry{left=1.15cm,right=1.15cm,top=1.1cm,bottom=1.1cm,columnsep=1.0cm}
\usepackage{paracol}
\usepackage{xeCJK}
\setmainfont{Roboto}
\setsansfont{Roboto}
\renewcommand{\familydefault}{\sfdefault}

% colors
\definecolor{SlateGrey}{HTML}{2E2E2E}
\definecolor{LightGrey}{HTML}{666666}
\definecolor{DarkPastelRed}{HTML}{450808}
\definecolor{PastelRed}{HTML}{8F0D0D}
\definecolor{GoldenEarth}{HTML}{E7D192}
\colorlet{name}{black}
\colorlet{tagline}{PastelRed}
\colorlet{heading}{DarkPastelRed}
\colorlet{headingrule}{GoldenEarth}
\colorlet{subheading}{PastelRed}
\colorlet{accent}{PastelRed}
\colorlet{emphasis}{SlateGrey}
\colorlet{body}{LightGrey}

% Change some fonts, if necessary
\renewcommand{\namefont}{\Huge\rmfamily\bfseries}
\renewcommand{\personalinfofont}{\footnotesize}
\renewcommand{\cvsectionfont}{\Large\rmfamily\bfseries}
\renewcommand{\cvsubsectionfont}{\large\bfseries}


% Change the bullets for itemize and rating marker
% for \cvskill if you want to
\renewcommand{\cvItemMarker}{{\small\textbullet}}
\renewcommand{\cvRatingMarker}{\faCircle}
\setlist{leftmargin=*,labelsep=0.45em,nosep,itemsep=0.12\baselineskip,topsep=0pt,after=\vspace{0.08\baselineskip}}
\renewcommand{\divider}{\textcolor{body!30}{\hdashrule{\linewidth}{0.6pt}{0.5ex}}\smallskip}
\renewcommand{\cvevent}[4]{%
  {\normalsize\color{emphasis}#1\par}
  \vspace{0.15ex}
  \ifstrequal{#2}{}{}{
  \textbf{\color{accent}#2}\par
  \vspace{0.15ex}}
  \ifstrequal{#3}{}{}{%
    {\small\makebox[0.5\linewidth][l]%
      {\BeginAccSupp{method=pdfstringdef,ActualText={\datename:}}\cvDateMarker\EndAccSupp{}%
      ~#3}%
    }}%
  \ifstrequal{#4}{}{}{%
    {\small\makebox[0.5\linewidth][l]%
      {\BeginAccSupp{method=pdfstringdef,ActualText={\locationname:}}\cvLocationMarker\EndAccSupp{}%
        ~#4}%
    }}\par
  \smallskip\normalsize
}


\begin{document}
\name{唐权威}
\tagline{Multimodal LLM / Speech-Text Retrieval / Agent Dev}
%% You can add multiple photos on the left or right
\photoR{2.8cm}{photo}
% \photoL{2.5cm}{Yacht_High,Suitcase_High}

\personalinfo{%
  % Not all of these are required!
  \email{1076451802@qq.com}
  \phone{17790556197}
  \homepage{blog.quanwei.fun}
  \github{tangquanwei}
  \printinfo{}{}

}
\mynames{Quanwei/Tang}

\makecvheader
%% Keep dense project bullets readable on one page.
\AtBeginEnvironment{itemize}{\small}

%% Set the left/right column width ratio to 6:4.
\columnratio{0.7}

% Start a 2-column paracol. Both the left and right columns will automatically
% break across pages if things get too long.
\begin{paracol}{2}
\cvsection{实习}

\cvevent{会议软件攻关 RD}{AISpeech 思必驰}{2025/04 - 2025/09}{苏州}
\begin{itemize}
\item \textbf{时间戳大模型训练}~ 实现"LLM生成时间戳，精确召回提升10\%+
\item \textbf{热词大模型训练}~ 多模态知识融合：结合语音和文本知识，降低语音识别错误率，提升语音问答准确率。
\item \textbf{检索推理大模型训练}~ 可解释性增强：输出外部知识引用路径，确保知识引用可溯源。
\end{itemize}

\divider

\cvevent{APA System / Agent RD}{Momenta 魔门塔}{2025/09 -- 2026/01}{苏州}
\begin{itemize}
\item \textbf{Agent开发} ~ 设计并开发任务流自动化 Agent，实现从任务提交提交、状态追踪到结果反馈的闭环处理，大幅降低跨平台操作的心智负担与响应延迟。
\item \textbf{自动FIT-Checker开发} ~ 构建高度自动化的开发工作流，输入需求文档，输出开发文档和PR。Agent根据需求文档查资料、写文档、写代码、测试、提交代码，人工只需要审核代码合入主线。开发了一系列 Fit-Checker 平均 P\&R >90\%，大大减少了人工检查的时间。
\end{itemize}

\divider

\cvevent{语音-热词跨模态检索 RD}{AISpeech 思必驰}{2026/01 - 2026/06}{苏州}
\begin{itemize}
\item \textbf{语音跨模态检索}：面向“语音中真实出现热词识别”任务，研发语音-热词检索方法，在语音流和文本知识之间直接建模，绕开 ASR 文本匹配的误差传递。
\item \textbf{Acoustic Token 对齐}：引入 CIF 将连续语音帧压缩为 acoustic tokens，设计长度感知滑窗与 token-level Text MaxSim，实现语音片段和文本 token 的细粒度匹配。
\item \textbf{训练与误差分析}：结合局部对比学习、CIF quantity 约束和辅助对齐目标训练；优化 batch 构造以过滤 false negative，通过 bad case 定位召回/误报来源。
\end{itemize}


\cvsection{项目}


\cvevent{基于大语言模型的医药问答 RAG 研究}{本科毕设}{2023/12 -- 2024/05}{}
\begin{itemize}
    \item \textbf{领域模型训练}：使用 LoRA/QLoRA 微调 ChatGLM-6B，完成医药 QA 领域知识注入，准确率较基线提升 10\%。
    \item \textbf{检索增强生成}：清洗 100k+ 原始数据得到 15,000 组 QA，使用 BGE-M3 向量化并结合 FAISS 索引，Top-5 召回率达到 98\%。
    \item \textbf{低延迟部署}：通过 4-bit 量化压缩模型显存占用，支持 CPU 端 2s 内端到端响应，关注生成内容的精炼性和可用性。
\end{itemize}

\divider

\cvevent{声学-文本跨模态对齐}{}{2026/01 -- 2026/06}{}
\begin{itemize}
    \item \textbf{跨模态检索建模}：面向语音热词检索任务，构建不依赖完整 ASR 转写的声学-文本跨模态检索框架；基于 Paraformer + BERT 双塔结构，实现大规模候选热词检索。
    \item \textbf{Acoustic Token 对齐}：引入 CIF 将帧级语音表示压缩为 acoustic tokens，通过 quantity loss 约束声学 token 与文本 token 的长度一致性；
    结合局部滑窗与 token-level MaxSim 定位长语音中的短热词片段。
    \item \textbf{两阶段训练与诊断}：实现 Stage1 CIF 预训练与 Stage2 联合检索训练，融合 local contrastive loss、CIF loss 和 token-level 辅助对齐目标，
    并处理 batch 内热词冲突导致的 false negative 问题。
    \item \textbf{部署与效果}：跑通数据处理、模型训练、量化部署和评估测试全流程；实现中文测试集 F1 96\%+，英文测试集 F1 85\%+，30s 语音、2k 热词检索延迟 $<$200ms。
\end{itemize}


\switchcolumn

\cvsection{教育}

\cvevent{苏州大学}{NLP Lab ~ 硕士}{2024/09 - 2027/06}{}

\divider

\cvevent{兰州交通大学}{软件工程 ~ 本科}{2020/09 - 2024/06}{}
  
\begin{itemize}
\item \textbf{专业排名：} 1/76 ~ \textbf{GPA：}3.89/4
\item \textbf{CTE-4:} 558 ~ \textbf{CET-6:} 478   
\end{itemize}

\cvsection{论文}

\begin{itemize}
    \item \textbf{ACL 2026 Findings（一作）}：\textit{Don't Just Listen, Try Planning: Graph-based Retrieval-Generation Agent for Long-form Audio Meeting Understanding}.
    \item \textbf{ACL 2025 Findings（一作）}：\textit{A Comprehensive Graph Framework for Question Answering with Mode-Seeking Preference Alignment}.
    \item \textbf{ICASSP 2026（一作）}：\textit{Relative Time Intervals Representation for Word-level Timestamping with Masked Training}.
    \item \textbf{ICASSP 2026（合作）}：\textit{TASU: Text-Only Alignment for Speech Understanding}.
\end{itemize}

\cvsection{技能}

\cvtag{Python}
\cvtag{PyTorch}
\cvtag{Deepspeed}
\cvtag{Multimodal LLM}
\cvtag{SFT/DPO}
\cvtag{Audio-Text Alignment}
\cvtag{Contrastive Learning}
\cvtag{LoRA/QLoRA}
\cvtag{CIF}
\cvtag{RAG}
\cvtag{FAISS/Milvus}
\cvtag{Agent}
\cvtag{Onnx}
\cvtag{TensorRT}
\cvtag{Neo4j}
\cvtag{Git}
\cvtag{Linux}
\cvtag{Docker}
\cvtag{Error Analysis}
% \cvtag{Pandas}
% \cvtag{Matplotlib}
% \cvtag{PostgreSQL/PGVector}
% \cvtag{FastAPI}
% \cvtag{NumPy}
% \cvtag{Java}
% \cvtag{Spring}
% \cvtag{JPA}

% \divider\smallskip

% \cvtag{Linux}
% \cvtag{Docker}
% \cvtag{VibeCoding}
% \cvtag{Neo4j}

\end{paracol}


\end{document}
